\section{Identification with a simple adjoint variety}
\label{hom:section}

We finish with the algebraic consequence of contact transitivity. The
argument uses the Borel fixed-point Theorem and standard highest-weight
theory, rather than a classification of homogeneous contact manifolds.

\begin{proposition}\label{hom:adjoint}
Let $(X,F)$ be a smooth connected complex projective contact manifold
with ample contact line $L$. Suppose a connected algebraic group $G$
acts effectively and transitively on $X$, preserving $F$. Then
$\mathfrak g=\Lie(G)$ is simple, and contact values of fundamental
vector fields define an isomorphism
\[
 \Phi:(X,F)\xrightarrow{\ \sim\ }
 \bigl(\PP(\mathcal O_{\min}(\mathfrak g)),F_{\mathrm{can}}\bigr).
\]
Under this isomorphism, $L$ is the restriction of the hyperplane line
bundle of $\PP(\mathfrak g)$.
\end{proposition}

\begin{proof}
The contact determinant formula gives $-K_X=L^{\otimes(n+1)}$,
so $X$ is Fano. The algebraic action canonically linearizes a sufficiently
high anticanonical power, giving an algebraic homomorphism to the projective
linear group of the resulting embedding. It has trivial kernel by
effectivity and, in characteristic zero, identifies $G$ with its closed
image. Thus $G$ is linear algebraic. Let
$R(G)$ be its connected solvable radical. The Borel fixed-point
Theorem for a connected solvable linear algebraic group acting on a
nonempty complete variety gives an $R(G)$-fixed point $x_0$
\cite[Corollary~17.3]{Milne2022}. Normality gives, for
$r\in R(G)$ and $g\in G$,
\[
 r(gx_0)=g(g^{-1}rg)x_0=gx_0.
\]
Transitivity and effectivity therefore imply $R(G)=1$. Thus $G$ is
semisimple.

For $D\in\mathfrak g$, write
$D_X(x)=\left.\frac{d}{dt}\right|_{0}\exp(tD)x$.
The quotient $L=T_X/F$ has its natural $G$-linearization. Transitivity
in characteristic zero makes evaluation of fundamental fields
surjective onto $T_X$, so their contact values give a surjection
\[
 \mathfrak g\otimes\OO_X\longrightarrow L,
 \qquad D\longmapsto\theta(D_X).
\]
Its dual is a line subbundle of $\mathfrak g^*\otimes\OO_X$.
With our convention that projective spaces parametrize lines, this
defines the $G$-equivariant morphism
\begin{equation}\label{hom:moment}
 \Phi:X\longrightarrow\PP(\mathfrak g^*),\qquad
 \Phi^*\OO_{\PP(\mathfrak g^*)}(1)\simeq L.
\end{equation}
It is finite onto its image: otherwise a positive-dimensional
projective fiber would contain a curve on which $L$ has degree zero,
contrary to ampleness. Its image is a single closed $G$-orbit, since
$X$ is projective and $G$ acts transitively.

Choose a Borel subgroup $B=TU$ of $G$ and a $B$-fixed point of $X$.
Its image is a $B$-stable line $\CC\ell\subset\mathfrak g^*$.
We spell out the highest-weight argument, including the possibility
of several simple factors. Decompose into simple ideals and their
dual representations:
\[
 \mathfrak g=\bigoplus_{i=1}^s\mathfrak g_i,
 \qquad \mathfrak g^*=\bigoplus_{i=1}^s\mathfrak g_i^*.
\]
Each $\mathfrak g_i^*$ is irreducible. Its $U$-fixed subspace is its
one-dimensional highest-weight space
\cite[Definition~22.16 and Theorem~22.18]{Milne2022}. The weights $\lambda_i$ of
these spaces are nonzero and distinct: on $\Lie(T)$ each is
supported on the Cartan subalgebra of the $i$th simple ideal, and is
zero on the other summands. These are the standard highest-weight
properties of the adjoint representation; see the proof of
\cite[Proposition~2.6]{Beauville1998} for the simple case.
Because a unipotent group has no nontrivial characters, $U$ fixes
$\ell$, while $T$ acts on $\CC\ell$ by a single character.
Writing $\ell=\sum_i\ell_i$, every nonzero $\ell_i$ would therefore
have that same weight $\lambda_i$. Distinctness forces exactly one
nonzero component. Thus $\CC\ell$ is the highest-weight line of one
simple coadjoint summand.

Let $x$ be the chosen point and put
$H=\operatorname{Stab}_G(x)$ and
$Q=\operatorname{Stab}_G([\ell])$. Equivariance gives $H\subset Q$,
and finiteness of $\Phi$ gives $\dim H=\dim Q$. The subgroup $Q$
contains $B$, so it is parabolic and is connected by
\cite[Corollary~17.49]{Milne2022}. This statement applies to every
connected semisimple $G$, without a simply connectedness hypothesis.
A connected algebraic group is irreducible;
its closed subgroup $H$ of the same dimension must therefore be
all of $Q$. Consequently the homogeneous map
\[
 X\simeq G/H\longrightarrow G/Q=\Phi(X)
\]
is an isomorphism. Every other simple factor acts trivially on the
chosen coadjoint summand, hence on $X$. Effectivity excludes these
factors, and $\mathfrak g$ is simple.

The Killing form now identifies $\mathfrak g^*$ with $\mathfrak g$
equivariantly. The highest-weight line becomes the line of a
highest-root vector $e$, which is nilpotent. Its orbit
$\mathcal O=G\cdot e$ is stable under nonzero scalar multiplication:
the nonzero root character of $T$ takes all values in $\CC^*$.
Thus $\PP(\mathcal O)=\Phi(X)$. For completeness, this is exactly
the projectivization of the unique smallest nonzero nilpotent orbit.
Indeed every nonzero nilpotent orbit is a cone: the
Jacobson--Morozov Theorem supplies an $\mathfrak{sl}_2$-triple and hence
the scalings. Its projective
closure contains a Borel-fixed point, necessarily the unique
highest-weight line. It therefore contains $\PP(\mathcal O)$.
Its dimension is at least $\dim\PP(\mathcal O)$, with equality
only when its projective closure equals that closed orbit. Since
both nilpotent orbits are cones, equality then gives the same
orbit. This also proves uniqueness; compare
\cite[Section~2.4 and Proposition~2.6]{Beauville1998}.

It remains to identify the given distribution. We do this on the
coadjoint cone, retaining the convention
\[
 (\ad^*(A)\ell)(D)=-\ell([A,D]),\qquad
 \omega_\ell(\ad^*(A)\ell,\ad^*(D)\ell)=\ell([A,D])
\]
for the coadjoint action and the Kirillov--Kostant--Souriau form.
The latter is a closed nondegenerate form on the orbit. Under
scalar multiplication it has weight one. If $E$ is the Euler
vector field and $\lambda=\iota_E\omega$, then
$d\lambda=\omega$ and $\lambda(E)=0$. The form $\lambda$ has weight
one, so it descends under $\mathcal O\to\PP(\mathcal O)$ to an
$\OO_{\PP(\mathfrak g^*)}(1)$-valued one-form. Explicitly, the cone is
the tautological line bundle with its zero section removed. For a tangent
vector $u$ at $[\ell]$, assign to $\ell$ the scalar
$\lambda_\ell(\widetilde u)$, where $\widetilde u$ is a lift of $u$.
The equality $\lambda(E)=0$ makes this independent of the lift,
and weight one makes it linear in $\ell$. Its value thus lies in
the dual tautological line. Its kernel is the canonical contact
distribution; contact nondegeneracy follows from
\[
 \lambda\wedge\omega^n
   =\frac{1}{n+1}\,\iota_E(\omega^{n+1})\ne0.
\]
This is the symplectic-cone construction of
\cite[Sections~2.1--2.2]{Beauville1998}.

At $\Phi(x)=[\ell]$, choose the nonzero functional on $L_x$ that
represents $\ell$ in the dual of the contact-value map. For
$\mathfrak h=\Lie(H)$, evaluation identifies $T_xX$ with
$\mathfrak g/\mathfrak h$, and
\begin{equation}\label{hom:hyperplane}
 \mathfrak h\subset\ker\ell,\qquad
 F_x=(\ker\ell)/\mathfrak h.
\end{equation}
Choose $A_0\in\mathfrak g$ with $\ad^*(A_0)\ell=\ell$; at a
highest-weight vector one takes a Cartan element of weight one,
and conjugation gives the assertion throughout the orbit. This
represents the Euler vector, and our conventions give
\[
 \lambda_\ell(\ad^*(D)\ell)
 =\omega_\ell(\ad^*(A_0)\ell,\ad^*(D)\ell)
 =\ell([A_0,D])=-\ell(D).
\]
Its kernel on the projective tangent space is precisely
\eqref{hom:hyperplane}. Changing the sign convention for $\omega$
changes neither kernel. Thus $d\Phi$ carries the original $F$ to
$F_{\mathrm{can}}$, as required; see also
\cite[Remark~2.3]{Beauville1998}.
\end{proof}

\begin{proof}[Completion of the proof of Theorem~\ref{thm:main}]
In the primitive Picard case, Corollary~\ref{jet:transitive} gives
connected algebraic contact transitivity. In the two other cases,
Proposition~\ref{red:alternatives} gives it directly. Apply
Proposition~\ref{hom:adjoint} to the effective identity component
of the contact automorphism group. The resulting algebraic
isomorphism identifies both $F$ and its quotient $L$, proving
Theorem~\ref{thm:main} for every $n\geq1$.
\end{proof}

\subsection{Consequences}

\begin{corollary}\label{hom:models}
The incidence and projective-space alternatives in
Proposition~\ref{red:alternatives} are the adjoint contact varieties
of $\mathfrak{sl}_{n+2}(\CC)$ and
$\mathfrak{sp}_{2n+2}(\CC)$, respectively. Their adjoint
polarizations are $\OO(1,1)$ and $\OO(2)$.
\end{corollary}

\begin{proof}
For $V=\CC^{n+2}$ the incidence embedding is
\[
 ([v],[f])\longmapsto[v\otimes f]\in\PP(\mathfrak{sl}(V)),
 \qquad f(v)=0.
\]
These are the rank-one nilpotent matrices, the orbit of a
highest-root line. The embedding is the restricted Segre embedding,
so its hyperplane line is $\OO(1,1)$. For the tautological contact
form the contact value of $D\in\mathfrak{sl}(V)$ is
$f(Dv)=\operatorname{tr}((v\otimes f)D)$; the Killing form is a
nonzero multiple of this trace pairing. This is therefore precisely
the projective contact-value map of Proposition~\ref{hom:adjoint}.

For a symplectic space $(W,\sigma)$ of dimension $2n+2$, put
$N_v(w)=\sigma(v,w)v$. Then $N_v\in\mathfrak{sp}(W)$, and
\[
 [v]\longmapsto[N_v]\in\PP(\mathfrak{sp}(W))
\]
is the quadratic Veronese embedding under
$\mathfrak{sp}(W)\simeq\Sym^2W$. Its image is the orbit of the
highest-root line, and its hyperplane line is $\OO_{\PP(W)}(2)$.
The symplectic projective contact form has value
$\sigma(v,Dv)=\operatorname{tr}(N_vD)$ on a fundamental field,
so this too is the contact-value map under the Killing pairing.
Both identifications preserve the given contact distributions by
Propositions~\ref{red:alternatives} and~\ref{hom:adjoint}; compare
\cite[Examples~2.7]{Beauville1998}.
\end{proof}

\begin{corollary}\label{jet:good-curves}
In the primitive Picard case, every point of $X$ lies on an
immersed rational curve $u:\PP^1\to X$ satisfying
\[
u^*L\simeq\OO(2),\qquad
\theta\,du:T_{\PP^1}\xrightarrow{\sim}u^*L,
\qquad u^*F\simeq\OO(1)^{\oplus2n}.
\]
\end{corollary}

\begin{proof}
Lemma~\ref{jet:value-velocity} gives a map $u$ near $g$ with
$\theta(u'(0))\ne0$. We may keep $u^*T_X$ strictly positive:
this is the open condition $H^1(u^*T_X(-2))=0$.
Its rank is $d=2n+1$ and its degree is
$(-K_X)\cdot u_*(\PP^1)=2(n+1)=d+1$. Therefore its splitting is
\[
u^*T_X\simeq\OO(2)\oplus\OO(1)^{\oplus2n}.
\]
The map $\theta\,du:\OO(2)\to\OO(2)$ is nonzero at $0$, hence
an isomorphism. In particular $u$ is immersive. Its differential
maps into the unique degree-two summand, because
$\Hom(\OO(2),\OO(1))=0$, and identifies that summand with
$T_{\PP^1}$. The contact quotient therefore splits by $du$, and
$u^*F\simeq\OO(1)^{\oplus2n}$. Finally,
Corollary~\ref{jet:transitive} carries $u(0)$ to any prescribed
point by a contact automorphism, preserving all these properties.
\end{proof}

\begin{proof}[Proof of Corollary~\ref{thm:contact-fano-ke}: K\"ahler--Einstein existence]
Theorem~\ref{thm:main} gives an algebraic isomorphism
$\varphi:X\to Y_{\mathfrak g}$. Wolf's correspondence identifies
$Y_{\mathfrak g}$, with its complex structure, as the twistor space of
a compact symmetric positive quaternionic-K\"ahler model
\cite{Wolf1965}; see also \cite[Section~2]{LeBrun1995}.
A constant rescaling normalizes the model metric without changing its
twistor complex structure. LeBrun's positive twistor theorem then gives
a K\"ahler--Einstein metric $h$ on $Y_{\mathfrak g}$ with
$\operatorname{Ric}_h=\lambda h$ for some $\lambda>0$; the theorem
includes $n=1$ under the four-dimensional self-dual Einstein convention
\cite[Theorem~2.1 and Definition~2.3]{LeBrun1995}.
Since $\varphi$ is biholomorphic, $\varphi^*h$ is K\"ahler for the
given complex structure on $X$, and naturality of Ricci curvature gives
$\operatorname{Ric}_{\varphi^*h}=\lambda\varphi^*h$.
\end{proof}

\begin{proof}[Proof of Corollary~\ref{thm:projective-contact}: projective contact manifolds]
An ample divisor has a nonzero class in $H^2(X,\mathbb Q)$, so
$b_2(X)\geq1$. Projectivity also makes $X$ compact K\"ahler. If
$b_2(X)=1$, Demailly's Corollary~3~\cite{Demailly2000} gives ampleness
of $-K_X$, and Theorem~\ref{thm:main} identifies the given contact
distribution with the canonical one on an adjoint variety.

If $b_2(X)\geq2$, Demailly's Corollary~4~\cite{Demailly2000}, obtained
from Theorem~1.1 of Kebekus--Peternell--Sommese--Wi\'sniewski
\cite{KPSW2000}, gives the asserted cotangent-bundle isomorphism over
a smooth projective manifold $Z$. The base is connected because it is
the image of $X$ under the bundle projection, and
$\dim X=2\dim Z-1$ gives $\dim Z=n+1$.

The adjoint varieties carry the canonical contact structures recalled
above. For the cotangent bundles, let $\pi:\PP(T^*Z)\to Z$ be the
projection. At $(z,[\xi])$, the condition
$\xi(d\pi(v))=0$ on a tangent vector $v$ defines the tautological contact
distribution, with contact line $\OO_{\PP(T^*Z)}(1)$
\cite[Section~1]{Demailly2000}. This proves the converse.
\end{proof}
